
\begingroup
\let\clearpage\relax
\chapter{METODOLOGIA}\label{sec:metodologia}
\endgroup


A metodologia descreve os procedimentos e técnicas utilizados para conduzir a pesquisa. Esta seção deve apresentar, de forma clara e detalhada, os métodos adotados, os instrumentos de coleta de dados, a população ou amostra estudada, e as técnicas de análise. A escolha da metodologia deve ser justificada, demonstrando sua adequação aos objetivos da pesquisa.

\section{Tipo de Pesquisa}
\label{sec:tipo_pesquisa}

Este trabalho adota uma abordagem [descreva o tipo de pesquisa, por exemplo: quantitativa, qualitativa ou mista]. A pesquisa [quantitativa/qualitativa/mista] foi escolhida porque [justifique brevemente a escolha]. Para \citeonline{lamport1986latex}, esse tipo de abordagem é adequado para [explique por que é adequado].

\section{Procedimentos Metodológicos}
\label{sec:procedimentos}

Os procedimentos metodológicos foram divididos em etapas principais, conforme descrito a seguir:

\begin{itemize}
    \item \textbf{Definição do problema} [Descreva como o problema de pesquisa foi definido].
    \item \textbf{Revisão da literatura} [Explique como foi realizada a revisão bibliográfica].
    \item \textbf{Coleta de dados} [Descreva os métodos e instrumentos utilizados para coletar dados].
    \item \textbf{Análise dos dados} [Explique as técnicas de análise utilizadas].
\end{itemize}

\section{População e Amostra}
\label{sec:populacao_amostra}

A população deste estudo consiste em [descreva a população estudada]. A amostra foi selecionada por meio de [descreva o método de amostragem, por exemplo: amostragem aleatória simples, estratificada, etc.], totalizando [número de participantes ou unidades de análise]. A escolha desse método de amostragem justifica-se por [explique a justificativa].

\section{Instrumentos de Coleta de Dados}
\label{sec:instrumentos}

Para a coleta de dados, foram utilizados os seguintes instrumentos:

\begin{itemize}
    \item \textbf{Instrumento 1} [Descreva o instrumento e sua aplicação].
    \item \textbf{Instrumento 2} [Descreva o instrumento e sua aplicação].
    \item \textbf{Instrumento 3} [Descreva o instrumento e sua aplicação].
\end{itemize}

Cada instrumento foi escolhido com base em [justifique a escolha dos instrumentos]. Para garantir a validade e confiabilidade dos dados, [descreva as medidas adotadas, como pré-testes ou validação por especialistas].

\section{Técnicas de Análise de Dados}
\label{sec:analise_dados}

Os dados coletados foram analisados por meio das seguintes técnicas:

\begin{itemize}
    \item \textbf{Técnica 1}Utilizou-se o software [nome do software], que permite [explique brevemente as funcionalidades do software].
    \item \textbf{Técnica 2} Adotou-se [descreva a técnica e sua aplicação].
\end{itemize}

\section{Limitações da Pesquisa}
\label{sec:limitacoes}

Este estudo apresenta algumas limitações, tais como:

\begin{itemize}
    \item \textbf{Limitação 1} [Explique como isso pode afetar os resultados].
    \item \textbf{Limitação 2} [Explique como isso pode afetar os resultados].
    \item \textbf{Limitação 3} [Explique como isso pode afetar os resultados].
\end{itemize}

Apesar dessas limitações, acredita-se que os resultados obtidos são válidos e relevantes para [contexto da pesquisa].

\section{Considerações Éticas}
\label{sec:etica}

Este estudo foi conduzido em conformidade com as diretrizes éticas para pesquisas envolvendo [descreva o tipo de participação, por exemplo: seres humanos, animais, etc.]. O projeto foi aprovado pelo Comitê de Ética em Pesquisa da [instituição], sob o número de protocolo [número do protocolo]. Todos os participantes assinaram um Termo de Consentimento Livre e Esclarecido (TCLE), garantindo o sigilo e a confidencialidade das informações coletadas.
